\documentclass[11pt]{article}
\usepackage[margin=0.7in]{geometry}
\usepackage{amsmath,amssymb}
\usepackage{enumitem}
\usepackage{tikz}
\usepackage{pgfplots}
\pgfplotsset{compat=1.18}
\usepackage{fancyhdr}
\usepackage{parskip}
\setlist[enumerate]{leftmargin=*, itemsep=6pt}
\pagestyle{fancy}
\fancyhf{}
\fancyfoot[L]{\small\textit{en}Vision\textsuperscript{\textregistered} Algebra 2}
\fancyfoot[C]{\small 6-2 Lesson Quiz, short form}
\fancyfoot[R]{\small\thepage}
\renewcommand{\headrulewidth}{0pt}
\newcommand{\choice}[1]{\tikz[baseline=-0.6ex]\node[draw, circle, inner sep=0.6pt, minimum size=3.3mm, font=\scriptsize\bfseries] {#1};\;}
\newcommand{\tile}[1]{\tikz[baseline=-0.6ex]\node[draw, rounded corners=2pt, minimum height=7mm, minimum width=14mm, inner sep=2pt] {$#1$};}
\newcommand{\expbox}{\tikz[baseline=-0.6ex]\node[draw, minimum height=5mm, minimum width=12mm] {};}
\newcommand{\ansbox}{\tikz[baseline=-0.6ex]\node[draw, minimum height=8mm, minimum width=18mm] {};}

\begin{document}

\noindent\textbf{Name}\; \rule{6cm}{0.4pt} \hfill Period \rule{1.2cm}{0.4pt} \quad Date \rule{2.2cm}{0.4pt}

\vspace{4pt}\noindent\rule{\textwidth}{0.6pt}

\vspace{2pt}\noindent{\LARGE\textbf{6-2 \; Lesson Quiz: Exponential Models}}

\noindent{\small short form \hfill Savvas numbers 3, 4, 5 (items 1 and 2: full Savvas quiz only); show your work}

\vspace{2pt}\noindent\rule{\textwidth}{0.6pt}

\vspace{6pt}
\textbf{3.} Micah invests \$5,280 in an account that earns 4.2\% interest, compounded monthly.

Complete the function that models the amount $A$ in the account after $t$ years. Write one tile in each box.

\begin{minipage}[c]{0.40\textwidth}
\centering
\tile{5280}\; \tile{5501}\; \tile{5298.48}\\[5pt]
\tile{\tfrac{0.042}{12}}\; \tile{\tfrac{0.042}{6}}\; \tile{0.042}\\[5pt]
\tile{12t}\; \tile{t}\; \tile{\tfrac{t}{12}}
\end{minipage}\hfill
\begin{minipage}[c]{0.56\textwidth}
\centering
\Large $A(t)=$ \ansbox $\Bigl(1+$ \ansbox $\Bigr)^{\raisebox{-2pt}{\expbox}}$
\end{minipage}

\vspace{2pt}
Work:

\vspace{0.5cm}

\noindent\rule{\textwidth}{0.4pt}

\vspace{6pt}
\textbf{4.} Steve invests \$1,800 in an account that earns 3.7\% annual interest, compounded continuously. What is the approximate value of the account after 10 years?

\vspace{2pt}
\choice{A} \$2,466 \hspace{1.2cm} \choice{B} \$2,589 \hspace{1.2cm} \choice{C} \$2,601 \hspace{1.2cm} \choice{D} \$2,606

Work:

\vspace{1.3cm}

\noindent\rule{\textwidth}{0.4pt}

\vspace{6pt}
\textbf{5.} Paul considers investing \$8,000 in three different accounts for 6 years.

\begin{minipage}[c]{0.55\textwidth}
Option A: 5\% simple interest\\
Option B: 4.5\% interest compounded continuously\\
Option C: 4.75\% interest compounded yearly
\end{minipage}\hfill
\begin{minipage}[c]{0.42\textwidth}
\fbox{\parbox{\dimexpr\linewidth-2\fboxsep-2\fboxrule}{\small\raggedright\textbf{Word bank:}\\ linear \quad exponential \\ quadratic \quad constant}}
\end{minipage}

\vspace{4pt}
After 6 years, Option B, to the nearest dollar, will be worth \$\rule{1.6cm}{0.4pt} more than Option A and \$\rule{1.6cm}{0.4pt} less than Option C.

\vspace{4pt}
Option A is a \rule{2.6cm}{0.4pt} function while Options B and C are \rule{2.6cm}{0.4pt} functions.

\vspace{4pt}
While the \rule{2.6cm}{0.4pt} function gains value the fastest initially, over time the \rule{2.6cm}{0.4pt} functions will gain value faster.

Work (use the back if you need more room):

\vfill
\noindent\textbf{Scoring:} Questions 3 and 4, 10 points each; Question 5, bonus up to 5 points. Total out of 20. (Question numbers are the Savvas quiz numbers; items 1 and 2 are on the full Savvas quiz only.)

\newpage
\noindent{\LARGE\textbf{6-2 Lesson Quiz, short form: answer key}}

\vspace{4pt}\noindent\rule{\textwidth}{0.6pt}

\begin{enumerate}
\item[\textbf{3.}] \textbf{$A(t)=5280\left(1+\frac{0.042}{12}\right)^{12t}$.} $P=5280$, $r=0.042$, $n=12$: $A=P\left(1+\frac rn\right)^{nt}$. Distractor tiles: $5501\approx5280(1.042)$ is one year at annual compounding; $5298.48=5280\left(1+\frac{0.042}{12}\right)$ is the value after one month; $\frac{0.042}{6}$ and $\frac t{12}$ mix up the number of periods. Savvas quiz item 3, DOK 2, Skills Review and Practice F67. Taught by Example 2 (Try It 2 is the same monthly model). \textbf{10 points:} 4 for $5280$, 3 for $\frac{0.042}{12}$, 3 for $12t$.
\item[\textbf{4.}] \textbf{D, \$2,606.} $A=Pe^{rt}=1800e^{0.037(10)}=1800e^{0.37}\approx1800(1.447735)\approx\$2{,}605.92$. Choice A is simple interest, $1800(1+0.37)=2466$; B is annual compounding, $1800(1.037)^{10}\approx2588.57$; C is quarterly compounding, $1800\left(1+\frac{0.037}{4}\right)^{40}\approx2601.49$. Savvas quiz item 4, DOK 2, F67. Taught by Example 4 (Try It 4). \textbf{10 points:} 6 for D, 4 for the work ($e$ raised to the whole product $0.037\cdot10$).
\item[\textbf{5.}] \textbf{80; 89; linear; exponential; linear; exponential.} A: $8000(1+0.05\cdot6)=\$10{,}400$. B: $8000e^{0.045(6)}=8000e^{0.27}\approx\$10{,}479.72$, so $10{,}480-10{,}400=80$. C: $8000(1.0475)^6\approx\$10{,}568.52$, so $10{,}569-10{,}480=89$. After one year A is $\$8{,}400$, B $\$8{,}368.22$, C $\$8{,}380$: the linear option gains the most at first; the exponential options gain faster over time. Savvas quiz item 5, DOK 2, F67. Taught by Examples 2 and 4 and rehearsed by the packet's Practice (Lena, \$5,000, 8 years: a structural parallel with new numbers). \textbf{Bonus 5:} 1 for each dollar amount (2); up to 3 for the four words, $\max(0,\,3-\text{number of wrong words})$.
\end{enumerate}

\textbf{Items left off this form} (full Savvas quiz only; available for a retake or Form B): Savvas item 1, the town of 4,200 in 1970 decreasing 2.4\% per year: $y=4200(0.976)^t$ and about $0.605\%$ per quarter ($1-0.976^{1/4}\approx0.006055$); DOK 2, F14; taught by Example 1, cut in the fast packet. Savvas item 2, the model $y=a\cdot b^x$ through $(3,5)$ and $(4,10)$: answer D, $(5,20)$ ($b=\frac{10}{5}=2$, so the next $y$ is $10\cdot2=20$); DOK 2, F13; taught by Example 5, cut. The quiz has no graphs, so nothing was left off for a figure.

\textbf{IXL replacement:} SmartScore 70 or above on YJW (Compound interest: word problems), 5GU (Continuously compounded interest: word problems) and KLF (Describe linear and exponential growth and decay) replaces this quiz with 20/20.

\end{document}
